% GENERATED FILE. Edit canonical lessons, then run npm run generate:books.
% canonical-source-hash: fnv1a64:ee651d55dc0eda18
% canonical-lessons: FR-C98-hier, FR-C98-souvent, FR-C98-parfois, FR-C98-ville, FR-C98-rue

\chapter{Yesterday, Often, Sometimes, Town, Street}
\label{ch:fr-yesterday-often-sometimes-town-street}
\hlchaptermodality{98}

\begin{chapteropening}
\textbf{By the end of this chapter:} \emph{I can name yesterday, often, sometimes, a town, a street.}

Complete the last lesson of chapter 98: I can name yesterday, often, sometimes, a town, a street.
\end{chapteropening}

\section[hier]{hier — yesterday}

\label{lesson:FR-C98-hier}



\begin{quote}
Before the new one: say the French for outside, then the French for today.
\end{quote}

\subsection*{You'll want to know: hier}
\textbf{hier} — \textquotedblleft{}yesterday\textquotedblright{}.

\begin{cousinweb}
\textbf{hier} $\leftarrow$ Old French \emph{ier} $\leftarrow$ Latin \emph{heri}, \textquotedblleft{}yesterday\textquotedblright{} — inherited. (EtymDB 2.1)
\end{cousinweb}

\begin{grammarlens}[title={What you've built}]
One more word for this chapter.
\end{grammarlens}

\subsection*{Your turn}
\begin{itemize}
\raggedright
  \item \emph{Say it:} \emph{hier}
  \item \emph{Say it:} \emph{hier}, once more
  \item \emph{Say it:} say \emph{hier}
  \item \emph{Recall:} say the French for outside, then the French for today, then say \emph{hier} again
\end{itemize}

\begin{tcolorbox}[breakable,colback=teal!4,colframe=teal!35!black,title={Before you move on}]
What does hier mean? (\textquotedblleft{}Yesterday\textquotedblright{}.) Say it once more.
\end{tcolorbox}
\section[souvent]{souvent — often}

\label{lesson:FR-C98-souvent}



\begin{quote}
Before the new one: say the French for today, then the French for yesterday.
\end{quote}

\subsection*{You'll want to know: souvent}
\textbf{souvent} — \textquotedblleft{}often\textquotedblright{}.

\begin{cousinweb}
\textbf{souvent} $\leftarrow$ Old French \emph{suvent} (EtymDB 2.1) $\leftarrow$ Latin \emph{subinde}, \textquotedblleft{}soon after, time after time\textquotedblright{} (Lewis \& Short) — inherited.
\end{cousinweb}

\begin{grammarlens}[title={What you've built}]
One more word for this chapter.
\end{grammarlens}

\subsection*{Your turn}
\begin{itemize}
\raggedright
  \item \emph{Say it:} \emph{souvent}
  \item \emph{Say it:} \emph{souvent}, once more
  \item \emph{Say it:} say \emph{souvent}
  \item \emph{Recall:} say the French for today, then the French for yesterday, then say \emph{souvent} again
\end{itemize}

\begin{tcolorbox}[breakable,colback=teal!4,colframe=teal!35!black,title={Before you move on}]
What does souvent mean? (\textquotedblleft{}Often\textquotedblright{}.) Say it once more.
\end{tcolorbox}
\section[parfois]{parfois — sometimes}

\label{lesson:FR-C98-parfois}



\begin{quote}
Before the new one: say the French for yesterday, then the French for often.
\end{quote}

\subsection*{You'll want to know: parfois}
\textbf{parfois} — \textquotedblleft{}sometimes\textquotedblright{}.

\begin{cousinweb}
\textbf{parfois} — built from the word for \textquotedblleft{}by\textquotedblright{} and the word for \textquotedblleft{}a time, an occasion\textquotedblright{}, $\leftarrow$ Latin \emph{vices}, \textquotedblleft{}turns\textquotedblright{} (Lewis \& Short): by times.
\end{cousinweb}

\begin{grammarlens}[title={What you've built}]
One more word for this chapter.
\end{grammarlens}

\subsection*{Your turn}
\begin{itemize}
\raggedright
  \item \emph{Say it:} \emph{parfois}
  \item \emph{Say it:} \emph{parfois}, once more
  \item \emph{Say it:} say \emph{parfois}
  \item \emph{Recall:} say the French for yesterday, then the French for often, then say \emph{parfois} again
\end{itemize}

\begin{tcolorbox}[breakable,colback=teal!4,colframe=teal!35!black,title={Before you move on}]
What does parfois mean? (\textquotedblleft{}Sometimes\textquotedblright{}.) Say it once more.
\end{tcolorbox}
\section[la ville]{la ville — a town, a city}

\label{lesson:FR-C98-ville}



\begin{quote}
Before the new one: say the French for often, then the French for sometimes.
\end{quote}

\subsection*{You'll want to know: la ville}
\textbf{la ville} — \textquotedblleft{}a town, a city\textquotedblright{}.

\begin{cousinweb}
\textbf{ville} $\leftarrow$ Latin \emph{villa}, \textquotedblleft{}a country house, a farm\textquotedblright{} — inherited. (EtymDB 2.1)
\end{cousinweb}

\begin{grammarlens}[title={What you've built}]
One more word for this chapter.
\end{grammarlens}

\subsection*{Your turn}
\begin{itemize}
\raggedright
  \item \emph{Say it:} \emph{la ville}
  \item \emph{Say it:} \emph{la ville}, once more
  \item \emph{Say it:} say \emph{la ville}
  \item \emph{Recall:} say the French for often, then the French for sometimes, then say \emph{la ville} again
\end{itemize}

\begin{tcolorbox}[breakable,colback=teal!4,colframe=teal!35!black,title={Before you move on}]
What does la ville mean? (\textquotedblleft{}A town, a city\textquotedblright{}.) Say it once more.
\end{tcolorbox}
\section[la rue]{la rue — a street}

\label{lesson:FR-C98-rue}



\begin{quote}
Before the new one: say the French for sometimes, then the French for a town.
\end{quote}

\subsection*{You'll want to know: la rue}
\textbf{la rue} — \textquotedblleft{}a street\textquotedblright{}.

\begin{cousinweb}
\textbf{rue} $\leftarrow$ Latin \emph{ruga}, \textquotedblleft{}a wrinkle, a furrow\textquotedblright{} (Lewis \& Short) — inherited: a street was a furrow between houses.
\end{cousinweb}

\begin{grammarlens}[title={What you've built}]
One more word for this chapter.
\end{grammarlens}

\subsection*{Your turn}
\begin{itemize}
\raggedright
  \item \emph{Say it:} \emph{la rue}
  \item \emph{Say it:} \emph{la rue}, once more
  \item \emph{Say it:} say \emph{la rue}
  \item \emph{Recall:} say the French for sometimes, then the French for a town, then say \emph{la rue} again
\end{itemize}

\begin{tcolorbox}[breakable,colback=teal!4,colframe=teal!35!black,title={Before you move on}]
What does la rue mean? (\textquotedblleft{}A street\textquotedblright{}.) Say it once more.
\end{tcolorbox}
